\documentclass[11pt,a4paper]{article}
\usepackage[utf8]{inputenc}
\usepackage[T1]{fontenc}
\usepackage{geometry}
\geometry{margin=1.8cm}
\usepackage{tcolorbox}
\usepackage{array}
\usepackage{tabularx}
\usepackage{lipsum}
\usepackage{pifont}
\usepackage{tikz}
\usepackage{graphicx}
\begin{document}

% --- En-tête élève ---
\begin{tcolorbox}[colback=white,colframe=black,title=Évaluation test %titre du chapitre/notion 
]
\textbf{Nom :} \hspace{4cm}
\textbf{Prénom :} \hspace{4cm}
\textbf{Classe :}
\end{tcolorbox}

% --- Autoévaluation  ---
\section*{Autoévaluation (à remplir par l’élève)}
\begin{tcolorbox}[colback=white,colframe=black]
\textbf{Qu’ai-je réussi ?}\\[1cm]
\textbf{Qu’est-ce qui a été difficile ? Pourquoi ?}\\[1.5cm]
\textbf{Que puis-je améliorer pour la prochaine évaluation ?}\\[1.5cm]
\end{tcolorbox}

% --- Titre de l'évaluation ---
\section*{Evaluation à remplir par l'enseignant}
\begin{tcolorbox}[colback=white,colframe=black]
Grille de compétences :
\begin{itemize}
    \item C1 — S’approprier
    \item C2 — Analyser/Raisonner.
    \item C3 — Réaliser.
    \item C4 — Valider
    \item C5 — Communiquer
\end{itemize}
\end{tcolorbox}
% --- Récapitulatif barème ---
\section*{Barème exercice/compétences}
\begin{tcolorbox}[colback=white,colframe=black,title=Critères de réussite]
\renewcommand{\arraystretch}{1.3}
\setlength{\extrarowheight}{0,3cm}
\begin{tabularx}{\linewidth}{|c|c|X|X|}
\hline
\textbf{N°Exercice} & \textbf{Points} & \textbf{Compétence(s) associée(s)}& Commentaires\\
\hline
Exercice 1 &  &
%Indiquer les compétences mobilisées dans l'exercice 
&
%indiquer les commentaires pour cet exercice
\\ 
\hline
Exercice 2 &  &%Indiquer les compétences mobilisées dans l'exercice 
&
%indiquer ici les commentaires pour cet exercice
\\
\hline
Exercice 3 &  &
%Indiquer les compétences mobilisées dans l'exercice 
&
%indiquer les commentaires pour cet exercice
\\
\hline
Exercice 4 &  & 
%Indiquer les compétences mobilisées dans l'exercice 
&
%indiquer les commentaires pour cet exercice
\\
\hline
\textbf{Note} &
\multicolumn{3}{c|}{\textbf{} 
%COMMENTAIRES a la main
}
\\
\hline
\end{tabularx}

\end{tcolorbox}

%-------------------------DEBUT EVAL--------------------------------------

ewpage
%coller ici les exos

    



\end{document}



%-----------QCM ----------------------------
\section*{N° d'exo}
\begin{enumerate}
    \item ECRIRE LA QUESTION ICI 
    
\tikz{\draw (0,0) rectangle (0.25,0.25);} A. %Ecrire ici la proposition

\tikz{\draw (0,0) rectangle (0.25,0.25);} B. %Ecrire ici la proposition

\tikz{\draw (0,0) rectangle (0.25,0.25);} C. %Ecrire ici la proposition

\tikz{\draw (0,0) rectangle (0.25,0.25);} D. %Ecrire ici la proposition

\end{enumerate}
%------------------Exercices ------------------
-> Ajouter un exo : 
\section*{Exercice NUMERO DE L'EXO : \textbf{%ECRIREenoncéICi
}}
\begin{enumerate}
    \item % AJOUTE QUESTION 
    \\[0.3cm]
    \item % AJOUTE QUESTION 
    \\[0.3cm]

\end{enumerate}

------------------------------------------------------
-> Ajouter Une question : à ajouter à la suite des \item AVANT le endenumerate :

  \item % AJOUTE QUESTION 
    \\[0.3cm]

-----------------------Ajouter de la place pour répondre : -------------------
Venir Augmenter la valeur du \\[0.3cm]
-----------------------Ajouter des lignes pointillets pour répondre : -------


oindent\dotfill

\vspace{0.5cm}


oindent\dotfill

\vspace{0.5cm}


oindent\dotfill



-> même chose : venir faire varier la valeur en cm pour adapter l'écart entre les lignes
-------------------------------------------------------------------------
Ajouter une image : 
-> comment faire : 
1 : Ajouter l'image en l'important en cliquant haut à gauche sur la petite fleche vers le haut
2: Renommer l'image
3: Coller le code ci dessous dans le fichier
4 : bien mettre le nom de l'image qui correspond au fichier importé dans les lignes de codes
\begin{figure}[h]
    \centering
    \includegraphics[width=0.5\textwidth]{NOMDELIMAGE.png}
    \caption{TITREDELIMAGESIBESOIN} %Si pas de titre : supprimer la ligne 
\end{figure}

-----Infos image : si besoin de faire varier la taille de l'image : le width=0.5 peut etre modifié : width=0.3 :  image plus petite, width 0,8 plus grande, width sans rien derrière : toute la largeur de la page
--------------------------------------------------

%-------------------Faire un exercice à relier ------------

\begin{tabular}{p{5cm} @{\hspace{5cm}} p{5cm}}
Chat \hfill {\Large$\bullet$}   & {\Large$\bullet$}\hspace{0.3cm} Abboie \\
Chien \hfill {\Large$\bullet$}  & {\Large$\bullet$}\hspace{0.3cm} Miaule \\
Oiseau \hfill {\Large$\bullet$} & {\Large$\bullet$}\hspace{0.3cm} Chante \\
\end{tabular}
